\section{Automated PD-L1 status prediction in lung cancer with multi-modal PET/CT fusion}
Da-Ano et al. (2024) developed a deep learning framework for non-invasive prediction of PD-L1 expression in non-small cell lung cancer (NSCLC) using multi-modal positron emission tomography/computed tomography (PET/CT) imaging \cite{Da-Ano2024-oo}. The study included 189 NSCLC patients from a single clinical center in France, each having both PET and CT scans acquired on Siemens Biograph mCT40 systems. The aim was to replace invasive biopsy-based immunohistochemistry with a fully automated image-based approach. Several convolutional neural network (CNN) architectures were evaluated, including ResNet-101, DenseNet-264, and EfficientNet-B1. Six data fusion strategies were tested to integrate PET and CT information: single-modality inputs (CT or PET), early fusion (combining both modalities at the input level), and late fusion (combining features at deeper network layers with or without shared weights).

All models were implemented in PyTorch and trained using cross-entropy loss with Adam and SGD optimizers on NVIDIA A6000 and RTX-series GPUs. The images were resampled and normalized before training, and the lesion regions were automatically segmented using a U-Net model trained on the public LIDC-IDRI dataset. Performance was assessed with area under the receiver operating characteristic curve (AUC) on five-fold cross-validation. The best results were obtained for the DenseNet-264 architecture with early fusion, achieving an AUC = 0.84 (95\% CI: 0.82-0.88), outperforming models trained on CT or PET alone.

The results demonstrated that combining anatomical (CT) and metabolic (PET) information enables more accurate PD-L1 status prediction than using either modality separately. Furthermore, early fusion proved more effective than late fusion in capturing complementary imaging features relevant for biomarker estimation. Overall, this study highlighted the potential of multimodal deep learning to support clinical decision-making by offering a non-invasive alternative to tissue-based PD-L1 testing.